\section{Related work}
\label{sec:related}

\paragraph{Temporal and bitemporal databases.}
Snodgrass and Ahn separated valid time, when a fact is true in the modelled
world, from transaction time, when it is stored~\cite{snodgrass1985taxonomy}; the
consensus glossary calls a database that records both
bitemporal~\cite{jensen1998glossary}, and Snodgrass's book gives the SQL
patterns for maintaining both by hand~\cite{snodgrass2000developing}.
SQL:2011 standardised application-time period tables, which carry a valid
interval as two columns, and system-versioned tables, which the database stamps
with transaction time; a table may be both~\cite{kulkarni2012sql2011}. The store
described here is bitemporal in that sense, with three differences. The
transaction coordinate is the later of a filer's claim and the server clock
rather than the server clock alone, so that a filer can record lateness but
cannot record earliness. Nothing is updated, so a transaction-time period has no
end column: a row is superseded by the existence of a later version, not by
being closed. And the unit of correction is a statement, a value from an
instant, rather than a row keyed by a primary key.

Datomic keeps an append-only log of facts, each stamped with the transaction
that asserted or retracted it, and its \texttt{as-of} and \texttt{since} views
filter by transaction~\cite{datomic}; valid time is modelled with ordinary
attributes. The first version of this store had the same shape, and
\S\ref{sec:eval-first} measures its consequence on history loaded after the
fact. XTDB is bitemporal by default: a write may state its valid time, which
otherwise defaults to the transaction time, and the system assigns the
transaction time~\cite{xtdb}. Helland argues for
append-only records from which current state is derived~\cite{helland2015immutability},
which is the design discipline here.

\paragraph{Statements about statements.}
RDF expresses provenance and time about a triple by reification, which describes
the triple with four further triples~\cite{rdf11semantics}; by named graphs,
which put the triple in a graph whose name carries the
provenance~\cite{carroll2005named}; by singleton properties, which mint a
property per statement~\cite{nguyen2014singleton}; or by RDF-star, which lets a
triple be the subject of another~\cite{hartig2017rdfstar}. Temporal RDF attaches
validity intervals to triples and defines a query semantics over
them~\cite{gutierrez2005temporal}. Wikidata attaches qualifiers such as start
time and end time to a statement and gives each statement a rank and
references~\cite{vrandecic2014wikidata}, and the choice among these encodings for
Wikidata has been compared directly~\cite{hernandez2015reifying}; YAGO2 attaches
time and space to its facts in a similar way~\cite{hoffart2013yago2}. Each of
these makes the statement addressable so that metadata can refer to it. Here the
statement is the row, and its provenance and interval are columns, so no
second-order statement is needed; the cost is that the provenance vocabulary is
fixed by the table rather than open.

\paragraph{Provenance.}
PROV-O models entities, activities and agents with relations such as
\texttt{wasGeneratedBy}, \texttt{wasAttributedTo} and
\texttt{wasDerivedFrom}, and an interval for an activity~\cite{lebo2013provo}.
The columns of an assertion map onto a small part of it: the filing identity
onto attribution, the source and evidence onto derivation, and the write onto a
generation time. The store records no activity and no chain of derivations, and
it exports nothing in PROV.

\paragraph{Temporal question answering over knowledge graphs.}
TempQuestions collected temporal questions over Freebase~\cite{jia2018tempquestions}.
CronQuestions provides a Wikidata-derived temporal graph and template questions
with annotations, together with CronKGQA, which answers them with temporal
knowledge-graph embeddings~\cite{saxena2021cronquestions}. TimeQA asks
time-sensitive questions over documents rather than a graph~\cite{chen2021timeqa}.
Our use of CronQuestions takes the annotation as given and measures the store
beneath a question-answering system, not the system; it is closer to a
correctness and cost test for temporal lookups than to a question-answering
result.

\paragraph{Temporal graphs for agent memory.}
Zep builds agent memory on Graphiti, a temporal knowledge graph whose edges carry
a validity interval in the world and a creation and expiry time in the system;
when a new edge contradicts an existing one, as judged by a language model, the
existing edge is invalidated by setting its end
instants~\cite{rasmussen2025zep,graphiti}. The two designs agree on keeping two
clocks. They differ in how an old fact stops holding: Graphiti writes end
instants onto the old edge, and decides contradiction with a model, while this
store leaves every row as written, derives the end of a value from later rows,
and decides what ends what by declared cardinality.

Semantica~\cite{semantica} is a Python library for knowledge and context graphs.
Its temporal model keeps \texttt{valid\_from} and \texttt{valid\_until} on each
relationship, and a \texttt{recorded\_at} and \texttt{superseded\_at} for the
transaction axis, and \texttt{query\_at\_time} takes an axis argument. At
0.7.0, \texttt{recorded\_at} is a field of the relationship that defaults to
\texttt{valid\_from} when absent; a caller who omits it therefore files history
that the transaction axis reports as known when it was so, which is the case
Proposition~\ref{prop:noback} excludes by deriving the instant on the server. The
lane measures its valid-time query only.

\paragraph{Communities and summaries.}
Louvain optimises modularity~\cite{newman2004modularity} by local moves and
aggregation~\cite{blondel2008louvain}; Leiden adds a refinement phase that
guarantees connected, and well-connected, communities~\cite{traag2019leiden}.
GraphRAG partitions an entity graph with Leiden, summarises each community with a
language model in advance, and answers a global question by mapping over the
summaries and reducing the partial answers~\cite{edge2024graphrag}. The answer
operation here keeps the map-reduce shape and changes three things: the partition
is computed from the rows at the requested point on each call, the map reads
labelled facts rather than stored summaries, and every citation in the result is
checked against what the model was given.

\paragraph{Erasure.}
Article~17 of the GDPR gives a data subject a right to erasure and exempts, among
other things, processing necessary for legal claims~\cite{gdpr}. This store makes
erasure the single exception to its append-only rule, lets a hold prevent it, and
records it by a digest rather than by content (\S\ref{sec:erasure}).
